\section{Introducción: panorama general del curso}

Esta clase es la séptima de MIT 6.S191 (Introducción al aprendizaje profundo) de la primavera de 2025, impartida por un profesor invitado de Google. El profesor dirige en Google el equipo de investigación aplicada de Gemini (Gemini Applied Research Group), encargado de desplegar los modelos Gemini en los distintos productos de Google, y también imparte un curso de aprendizaje automático en el programa de maestría en ciencia de datos de UC Berkeley. El contenido de esta clase proviene de su LLM Boot Camp interno en Google.

Esta clase cubre cuatro temas principales:
\begin{enumerate}
    \item \textbf{LLM 101}: los principios básicos y el mecanismo de funcionamiento de los modelos de lenguaje
    \item \textbf{Aplicaciones prácticas de los LLM}: cómo construir un chatbot a partir de un modelo de lenguaje
    \item \textbf{Prompt Engineering}: las técnicas centrales de la ingeniería de prompts
    \item \textbf{Beyond Prompts}: técnicas que van más allá del prompt: fine-tuning, alineación y agentes de IA
\end{enumerate}

El objetivo docente central del profesor es que los estudiantes entiendan cómo funcionan los LLM, que desarrollen una intuición sobre cuándo y cómo usarlos, y que conozcan sus trampas comunes (lo que la literatura llama la ``jagged frontier'', la frontera dispareja de capacidades).

